\ProvidesFile{section_8_svd}[Section 8]

\newpage

\section{Singular Value Decomposition}

\subsection{Motivation for SVD}

If $A$ is an $m \times n$ matrix of rank $r$, then it can be written as
$$
A = B C,
$$
where $B$ is an $m \times r$ matrix and  $C$ is an $r \times n$ matrix. Then instead of storing the whole matrix $A$ ($mn$ memory cells), we can store the matrices $B$ and $C$ in $(m+n)r$ memory cells.

\begin{exercise}
Consider the matrix
$$
A = \begin{pmatrix}
1 & 2 & -1 \\
3 & 6 & -3 \\
-1 & -2 & 1
\end{pmatrix}.
$$
Represent $A$ as $B C$, where $B$ is $3\times 1$ matrix and $C$ is $1\times 3$ matrix.
\end{exercise}

\begin{exercise}
Consider the matrix
$$
A = \begin{pmatrix}
1 & 2 & -1 \\
2 & 1 & 0 \\
4 & 5 & -2
\end{pmatrix}.
$$
Represent $A$ as $B C$, where $B$ is $3\times 2$ matrix and $C$ is $2\times 3$ matrix.
\end{exercise}

We will construct a representation $A = B D C$, where $D$ is a diagonal matrix.

Often the matrix $A$ has full rank; nevertheless, for compression and noise reduction we will later discard the smallest diagonal entries of $D$ (corresponding to the smallest singular values).

\subsection{Definition of SVD}

\begin{lemma}
Matrices $A A^{\mathrm T}$ and $A^{\mathrm T} A$ are symmetric and share the same non-zero eigenvalues.
\end{lemma}

Then there exist orthogonal matrices $U$ and $V$ and a diagonal matrix with nonnegative entries $\Lambda$ such that
$$
A A^{\mathrm T} = U \Lambda U^{\mathrm T},
$$
$$
A^{\mathrm T} A = V \Lambda V^{\mathrm T}
$$.

When $A$ is a real matrix, the symmetric matrices $A A^{\mathrm T}$ and $A^{\mathrm T} A$ have real eigenvalues. If $A$ has independent rows, then $A A^{\mathrm T}$ is positive definite and all its eigenvalues are positive. If it has independent columns, then $A^{\mathrm T} A$ is positive definite and all its eigenvalues are positive. In the general case, $A A^{\mathrm T}$ and $A^{\mathrm T} A$ are positive semidefinite and their eigenvalues are nonnegative.

Let us arrange the eigenvalues of $A A^{\mathrm T}$ (and $A^{\mathrm T} A$) in decreasing order:
$$
\sigma_1^2 \geqslant \sigma_2^2 \geqslant \cdots \geqslant \sigma_r^2 > 0 = \sigma_{r+1}^2 = \cdots = \sigma_n^2,
$$
where $r$ is the rank of $AA^T$ (and $A^TA$).

Denote a matrix $\Sigma = \diag(\sigma_1, \sigma_2, \cdots, \sigma_r, 0, \cdots, 0)$.

Then for orthogonal matrices $U$ (eigenvectors of $A A^{\mathrm T}$) and $V$ (eigenvectors of $A^{\mathrm T} A$) we can write
$$
AA^T = U \Sigma^2 U^{\mathrm T} = U\Sigma V^T V\Sigma U^T = (U\Sigma V^T)(U\Sigma V^T)^{\mathrm T}.
$$
We call $A = U \Sigma V^{\mathrm T}$ the Singular Value Decomposition (SVD) of $A$. The numbers $\sigma_1, \sigma_2, \ldots, \sigma_{\min(m,n)}$ are the singular values of $A$, and the columns of $U$ and $V$ are the left and right singular vectors of $A$, respectively. The number of nonzero singular values equals the rank of $A$.

(We skip the rigorous proof of the existence of SVD and move to the algorithm using the observations above.)

\subsection{Low-rank approximation}

\begin{definition}
The Frobenius norm of a matrix $A$ is defined as
$$
\|A\|_F = \sqrt{\sum_{i=1}^m \sum_{j=1}^n |a_{ij}|^2}.
$$
\end{definition}

\begin{theorem}
Fix a matrix $A$ and consider a problem
$$
\left\{\begin{array}{l}
    \|A-B\|_F \rightarrow \min _B \\
    \operatorname{rk} B \leqslant k
    \end{array}\right.
$$
Then the solution of this problem is the matrix $B = U \Sigma_k V^{\mathrm T}$, where $\Sigma_k$ is obtained from $\Sigma$ by keeping only the first $k$ singular values (and zeroing out the rest).
\end{theorem}

\subsection{Pseudoinverse}

\begin{exercise}
Let $A$ be a nondegenerate matrix with known singular values $\sigma_1, \sigma_2, \cdots, \sigma_n$. Find the singular values of $A^{-1}$.
\end{exercise}

\begin{definition}
The pseudoinverse of a matrix $A$ is defined as
$$
A^+ = V \Sigma^+ U^{\mathrm T},
$$
where $\Sigma^+$ is the matrix with the reciprocals of the non-zero singular values of $A$ on the diagonal.
\end{definition}

Note that if $A$ is invertible, then $A^+ = A^{-1}$.

Consider the following problem: solve the matrix equation $Ax = b$. If $A$ is invertible, then there is a unique solution $x = A^{-1}b$. If $A$ is not invertible, then there can be more than one solution or zero solutions. In case it does not have a solution, we can find the 'solution' $y$ that minimizes the norm of the error $\|Ay - b\|$. It turns out that the solution is given by $y = A^+b$.

\subsection{Problems}

\begin{problem}
Find the rank $r$ of matrix $A$ and represent it as a product of matrices with $r$ columns and $r$ rows:
\begin{enumerate}[a) ]
    \item 
    $
    A=\begin{pmatrix}
        1 & 1 & 1 & 1 & 1 & 1 \\
        1 & 1 & 1 & 1 & 1 & 1 \\
        1 & 1 & 1 & 1 & 1 & 1 \\
        1 & 1 & 1 & 1 & 1 & 1 \\
        1 & 1 & 1 & 1 & 1 & 1 \\
        1 & 1 & 1 & 1 & 1 & 1
        \end{pmatrix},
    $
    \item 
    $
    A=\begin{pmatrix}
        a & a & c & c & e & e \\
        a & a & c & c & e & e \\
        a & a & c & c & e & e \\
        a & a & c & c & e & e \\
        a & a & c & c & e & e \\
        a & a & c & c & e & e
        \end{pmatrix}.
    $
\end{enumerate}
\end{problem}

\begin{problem}
Find the SVD of the matrix
\begin{enumerate}[a) ]
\item
$\begin{pmatrix}
    1 & 0 \\
    1 & 1
    \end{pmatrix}$
\item
$\begin{pmatrix}
        1 & 0 & 0 & 0 \\
        1 & 1 & 0 & 0 \\
        1 & 1 & 1 & 0 \\
        1 & 1 & 1 & 1
        \end{pmatrix}$
\end{enumerate}
\end{problem}

\begin{problem}
Find eigenvalues, eigenvectors and the SVD of the matrix
\begin{enumerate}[a) ]
\item $\begin{pmatrix}
    3 & 0 \\
    4 & 5
    \end{pmatrix}$
\item $\begin{pmatrix}
        0 & 1 & 0 & 0 \\
        0 & 0 & 2 & 0 \\
        0 & 0 & 0 & 3 \\
        0 & 0 & 0 & 0
        \end{pmatrix}$
\item $\begin{pmatrix}
            0 & 4 \\
            0 & 0
            \end{pmatrix}$
\item $\begin{pmatrix}
            0 & 4 \\
            1 & 0
            \end{pmatrix}$
\item $\begin{pmatrix}
                2 & 2 \\
                -1 & 1
                \end{pmatrix}$
\item $\begin{pmatrix}
                    1 & 1 & 0 \\
                    0 & 1 & 1
                    \end{pmatrix}$
\end{enumerate}
\end{problem}

\begin{problem}
Find the pseudoinverse of the matrices from 8.3. Compare this to the inverse of the matrices which are invertible.
\end{problem}

\begin{problem}
Find the best approximation to the solution of the system of linear equations ($x$ is a column vector of corresponding size):
\begin{enumerate}[a) ]
    \item $\begin{pmatrix}
        0 & 1 & 0 & 0 \\
        0 & 0 & 2 & 0 \\
        0 & 0 & 0 & 3 \\
        0 & 0 & 0 & 0
        \end{pmatrix}x = \begin{pmatrix}
        0 \\ 1 \\ 2 \\ 3
        \end{pmatrix}$,
\item $\begin{pmatrix}
            0 & 4 \\
            0 & 0
            \end{pmatrix}x = \begin{pmatrix}
            1 \\ 1
            \end{pmatrix}$,
\item $\begin{pmatrix}
                2 & 2 \\
                -1 & 1
                \end{pmatrix}x = \begin{pmatrix}
                -1 \\ -1
                \end{pmatrix}$.
\end{enumerate}
\end{problem}